\unit{Exposé 232 -- Les schémas de Picard : théorèmes d'existence}
\sid{232}
\src{143}
\seminar{14e année, 1961/62, n\textsuperscript{o} 232\\Février 1962}
\paperhead{TECHNIQUE DE DESCENTE ET THÉORÈMES D'EXISTENCE EN GÉOMÉTRIE ALGÉBRIQUE\\
V. LES SCHÉMAS DE PICARD : THÉORÈMES D'EXISTENCE}%
  {par Alexander GROTHENDIECK}

\fsection{\normalfont 1}{\normalfont\uwave{Groupes et foncteurs de Picard relatifs}.}

Pour tout préschéma (plus généralement, tout espace annelé) $\symup{X}$, nous appelons
\uline{groupe de Picard} (absolu) de $\symup{X}$, et notons $\symup{\operatorname{Pic}(X)}$, le
groupe des classes, à un isomorphisme près, de Modules inversibles (i. e.
localement isomorphes à $\symup{\mathcal O_X}$) sur $\symup{X}$. On a donc un isomorphisme
canonique
\begin{equation}
\symup{\operatorname{Pic}(X)\overset{\sim}{\to}H^1(X,\mathcal O_X^*),
\tag{1.1}}
\end{equation}
où $\symup{\mathcal O_X^*}$ désigne le faisceau des unités de $\symup{\mathcal O_X}$, (qui
s'identifie en effet au faisceau des automorphismes du Module inversible type
$\symup{\mathcal O_X}$). Notons que $\symup{X\rightsquigarrow\operatorname{Pic}(X)}$ est un foncteur
contravariant en $\symup{X}$ de façon évidente, et que l'isomorphisme (1.1) est
fonctoriel.

Si $\symup{X}$ est un préschéma au-dessus d'un autre $\symup{S}$, alors pour un $\symup{S'}$ variable
dans la catégorie $\symup{(\text{\uwave{Sch}})/S}$ des préschémas sur $\symup{S}$, on a un foncteur
contravariant $\symup{S'\rightsquigarrow\operatorname{Pic}(X\times_S S')}$ grâce à ce qui
précède. Ce foncteur n'a aucune chance d'être \fquote{représentable} ([4], II, A),
car par suite de l'existence \uline{d'automorphismes non triviaux} des Modules
inversibles qu'on se propose de classifier, ce foncteur n'est pas de \fquote{\uline{nature locale}} ([5], IV, 5.4). Il y a donc lieu de le \fquote{rendre local}, en introduisant
de façon générale, pour tout préschéma relatif $\symup{X/S}$, un groupe de nature
relative,
\begin{equation}
\symup{\operatorname{Pic}'(X/S)=H^0\bigl(S,R^1f_*(\mathcal O_X^*)\bigr),
\tag{1.2}}
\end{equation}
(où $\symup{f:X\to S}$ est le morphisme structural) (comparer [4], II,
C 3). Dans loco citato, ce groupe est appelé groupe de Picard relatif, il sera
préférable ici de l'appeler \uline{groupe de Picard relatif restreint} de $\symup{X/S}$,
pour des raisons qui vont apparaître. Lorsque $\symup{S'}$ varie dans $\symup{(\text{\uwave{Sch}})/S}$,
$\symup{S'\rightsquigarrow\operatorname{Pic}'(X\times_S S'/S')}$ est un foncteur
contravariant en $\symup{S'}$, noté aussi $\symup{\mathop{\text{\uwave{Pic}}}\nolimits'_{X/S}}$, donc défini
essentiellement par la formule
\begin{equation}
\symup{\mathop{\text{\uwave{Pic}}}\nolimits'_{X/S}(S')=
\operatorname{Pic}'(X\times_S S'/S').
\tag{1.3}}
\end{equation}

\src{144}
Ce foncteur est maintenant \fquote{de nature locale}, vu qu'on a fait ce qu'il
fallait pour cela. Intuitivement, le deuxième membre de (1.3) s'interprète
comme l'ensemble des \fquote{familles algébriques} de classes de faisceaux
inversibles sur (les fibres de) $\symup{X/S}$, indexées par le préschéma de paramètres
$\symup{S'/S}$. Lorsque le foncteur $\symup{\mathop{\text{\uwave{Pic}}}\nolimits'}$ est représentable, le
préschéma sur $\symup{S}$ qui le représente est noté $\symup{\underline{\operatorname{Pic}}_{X/S}}$, et
appelé le \uline{préschéma de Picard} de $\symup{X}$ sur $\symup{S}$, donc, on aura alors
\begin{equation}
\symup{\operatorname{Hom}_S(S',\underline{\operatorname{Pic}}_{X/S})
\simeq\mathop{\text{\uwave{Pic}}}\nolimits'_{X/S}(S')
=\operatorname{Pic}'(X\times_S S'/S').
\tag{1.4}}
\end{equation}

Il y a cependant des cas importants où $\symup{\mathop{\text{\uwave{Pic}}}\nolimits'_{X/S}}$ n'est pas
représentable (exemple : variété de \fquote{Brauer-Severi} sur un corps $\symup{k}$, sans
point rationnel sur $\symup{k}$), et où cependant il existe une définition naturelle
d'un préschéma de Picard relatif. Cela tient au fait que dans la définition du
foncteur $\symup{\mathop{\text{\uwave{Pic}}}\nolimits'}$ à partir des groupes de Picard absolus
$\symup{\operatorname{Pic}(X\times_S S')}$, on n'a pas encore assez localisé, de
façon précise $\symup{\mathop{\text{\uwave{Pic}}}\nolimits'}$ n'est en général pas \fquote{compatible avec la descente fidèlement plate}. Explicitons.

Soit $\symup{(\mathfrak M)}$ l'ensemble des morphismes de préschémas qui sont
\uline{fidèlement plats et quasi-compacts}, cet ensemble est stable par
changement de base et par composition. Soit $\symup{P}$ un foncteur contravariant de
$\symup{(\text{\uwave{Sch}})/S}$ dans la catégorie des ensembles, et pour tout
$\symup{S}$-morphisme $\symup{u:T'\to T}$, $\symup{u\in(\mathfrak M)}$, considérons le
diagramme
\begin{equation}
\symup{P(T)\to P(T')\rightrightarrows P(T'\times_TT')
\tag{1.5}}
\end{equation}
transformé par $\symup{P}$ du diagramme
\[
\symup{T\leftarrow T'\mathrel{\overset{\mathrm{pr}_1,\ \mathrm{pr}_2}{\substack{\longleftarrow\\[-0.35ex]\longleftarrow}}}T'\times_TT'.}
\]
Si $\symup{P}$ est représentable, il résulte de la théorie de la descente ([4], I, B,
th. 2) que le diagramme (1.5) est exact pour tout $\symup{u\in(\mathfrak M)}$. On exprime
ce fait en disant que $\symup{P}$ est compatible avec $\symup{(\mathfrak M)}$, en l'occurrence
que $\symup{P}$ est \fquote{compatible avec la descente fidèlement plate}, ou encore que le
\fquote{\uline{préfaisceau}} $\symup{P}$ sur $\symup{(\text{\uwave{Sch}})/S}$ est un \fquote{\uline{faisceau}} pour la notion
de localisation fournie par l'ensemble $\symup{(\mathfrak M)}$. Lorsque $\symup{P}$ est
quelconque, un procédé standard, bien connu dans le cas de la localisation
topologique habituelle, permet de lui associer un \fquote{faisceau} $\symup{\mathcal P}$ et
un homomorphisme de foncteurs $\symup{P\to\mathcal P}$, qui soit universel
dans un sens évident. Le calcul de $\symup{\mathcal P}$ peut s'expliciter de la façon
suivante : pour définir $\symup{\mathcal P(T)}$, on désigne, pour tout $\symup{T'}$ sur $\symup{T}$ tel
que le morphisme $\symup{u:T'\to T\in(\mathfrak M)}$, par

\src{145}
$\symup{\overline H^0(T'/T,P)}$ le sous-ensemble de $\symup{P(T')}$ formé des éléments $\symup{\xi}$
dont les deux images $\symup{\xi_1,\xi_2}$ dans $\symup{P(T'\times_TT')}$ sont telles qu'il
existe un morphisme $\symup{v:T''\to T'\times_TT'}$, $\symup{v\in(\mathfrak M)}$,
tel que $\symup{\xi_1}$ et $\symup{\xi_2}$ aient même image dans $\symup{P(T'')}$.

[N. B. - L'ensemble
$\symup{\overline H^0}$ ainsi défini est donc plus grand que l'ensemble
$\symup{H^0(T'/T,P)}$ introduit dans [4], I, A, 4 (a).] Lorsque $\symup{T'}$ varie sur $\symup{T}$ fixé,
(toujours avec $\symup{u\in(\mathfrak M)}$) les $\symup{\overline H^0(T'/T,P)}$ forment un
diagramme (indexé par l'opposée de la catégorie des $\symup{T'}$ sur $\symup{T}$ et des
morphismes au-dessus de $\symup{T}$), et on pose
\begin{equation}
\symup{\mathcal P(T)=\varinjlim\limits_{T'}\overline H^0(T'/T,P).
\tag{1.6}}
\end{equation}
La loi fonctorielle en $\symup{T}$ de cette expression s'explicite de façon évidente.

Lorsque
\[
\symup{P(T)=\operatorname{Pic}(X\times_ST),}
\]
le foncteur contravariant sur $\symup{(\text{\uwave{Sch}})/S}$ défini par (1.6) est appelé
\uline{foncteur de Picard relatif} de $\symup{X}$ sur $\symup{S}$, et noté
$\symup{\mathop{\text{\uwave{Pic}}}\nolimits_{X/S}}$, et on appelle \uline{groupe de Picard relatif} de
$\symup{X}$ sur $\symup{S}$, et on note $\symup{\operatorname{Pic}(X/S)}$, le groupe
$\symup{\mathop{\text{\uwave{Pic}}}\nolimits_{X/S}(S)}$. On trouve alors une bijection évidente
\begin{equation}
\symup{\mathop{\text{\uwave{Pic}}}\nolimits_{X/S}(T)\overset{\sim}{\to}\operatorname{Pic}(X\times_ST/T).
\tag{1.7}}
\end{equation}

Un élément de $\symup{\operatorname{Pic}(X/S)}$ est donc défini à l'aide d'un élément
$\symup{\xi'}$ d'un groupe $\symup{\operatorname{Pic}(X\times_SS')}$ (où $\symup{S'\to S}$
est fidèlement plat et quasi-compact), tel que l'on puisse trouver un morphisme
fidèlement plat et quasi-compact $\symup{S''\to S'\times_SS'}$ tel que les
deux images inverses de $\symup{\xi'}$ dans $\symup{\operatorname{Pic}(X\times_SS'')}$ soient
les mêmes. Un élément $\symup{\xi'}$ de $\symup{\operatorname{Pic}(X\times_SS')}$ et $\symup{\xi_1}$
de $\symup{\operatorname{Pic}(X\times_SS_1)}$ (satisfaisant les conditions qu'on vient
d'expliciter), définissent le même élément de $\symup{\operatorname{Pic}(X/S)}$, si et
seulement s'il existe un morphisme fidèlement plat et quasi-compact
$\symup{S'_1\to S'\times_SS_1}$ tel que les images des deux éléments en
question dans $\symup{\operatorname{Pic}(X\times_SS'_1)}$ soient égales. Il est souvent
commode de travailler encore avec le foncteur
$\symup{P'=\mathop{\text{\uwave{Pic}}}\nolimits'_{X/S}}$ introduit plus haut, on constate aussitôt que
le morphisme canonique $\symup{P\to P'}$ définit un \uline{isomorphisme}
\begin{equation}
\symup{\mathcal P\overset{\sim}{\to}\mathcal P',
\tag{1.8}}
\end{equation}
ce qui donne une description de $\symup{\mathop{\text{\uwave{Pic}}}\nolimits_{X/S}}$ en termes de
$\symup{\mathop{\text{\uwave{Pic}}}\nolimits'_{X/S}=P'}$, plus commode souvent. En vertu de 2.3
ci-dessous, lorsque dans la description de $\symup{\operatorname{Pic}(X/S)}$ qu'on
vient d'expliciter, on remplace $\symup{P}$ par $\symup{P'}$, on peut en effet
\src{146}
faire $\symup{S''=S'\times_SS'}$, $\symup{S'_1=S'\times_SS_1}$, du moins sous les conditions
explicitées dans loco citato.

Lorsque le foncteur $\symup{\mathop{\text{\uwave{Pic}}}\nolimits_{X/S}}$ est représentable, on dit que
$\symup{X/S}$ admet un préschéma de Picard, et le préschéma sur $\symup{S}$ représentant le
foncteur est appelé \uline{préschéma de Picard} de $\symup{X}$ sur $\symup{S}$, et noté encore
$\symup{\underline{\operatorname{Pic}}_{X/S}}$. Il suffit évidemment pour ceci que
$\symup{P'=\mathop{\text{\uwave{Pic}}}\nolimits'_{X/S}}$ soit représentable, car alors $\symup{P'}$ est déjà un
``faisceau'', et la formule (1.8) prouve que le morphisme
$\symup{P'\to\mathcal P'}$ s'identifie au morphisme canonique
\begin{equation}
\symup{\mathop{\text{\uwave{Pic}}}\nolimits'_{X/S}\to\mathop{\text{\uwave{Pic}}}\nolimits_{X/S},
\tag{1.9}}
\end{equation}
qui est alors un isomorphisme. Donc notre terminologie est compatible avec
celle introduite plus haut avec (1.4). En général, lorsque
$\symup{\underline{\operatorname{Pic}}_{X/S}}$ existe, il est défini par l'isomorphisme fonctoriel :
\begin{equation}
\symup{\operatorname{Hom}_S(S',\underline{\operatorname{Pic}}_{X/S})\overset{\sim}{\to}
\operatorname{Pic}(X\times_SS'/S').
\tag{1.10}}
\end{equation}

\fsection{2}{\normalfont\uwave{Relations entre les divers groupes de Picard relatifs et absolus}.}

\result{PROPOSITION}{2.1} Soit $\symup{f:X\to S}$ un morphisme tel que
$\symup{\mathcal O_S\overset{\sim}{\to}f_*(\mathcal O_X)}$. Alors on a une suite exacte
\[
\symup{0\to\operatorname{Pic}(S)\to\operatorname{Pic}(X)
\to\operatorname{Pic}'(X/S).}
\]
Lorsque $\symup{X}$ admet une section sur $\symup{S}$, le dernier morphisme est surjectif, i. e.
on a un isomorphisme
\[
\symup{\operatorname{Pic}'(X/S)\simeq
\operatorname{Pic}(X)/\operatorname{Pic}(S).}
\]

La suite exacte peut être considérée comme la suite exacte en bas degrés
correspondant à la suite spectrale de Leray pour $\symup{f}$ et $\symup{\mathcal O_X^*}$. La
deuxième assertion également est formelle.

\result{PROPOSITION}{2.2} Soit $\symup{f:X\to S}$ un morphisme
quasi-compact et séparé tel que
$\symup{\mathcal O_S\overset{\sim}{\to}f_*(\mathcal O_X)}$, et soit
$\symup{S'\to S}$ un morphisme fidèlement plat et quasi-compact. Alors

\noindent (i) $\symup{\operatorname{Pic}'(X/S)\to
\operatorname{Pic}'(X\times_SS'/S')}$ est injectif ;

\noindent (ii) Si $\symup{X}$ admet une section localement sur $\symup{S}$ (i. e. tout $\symup{s\in S}$
a un voisinage ouvert $\symup{U}$ tel que $\symup{X|U}$ ait une section sur $\symup{U}$), alors le
diagramme
\[
\symup{\operatorname{Pic}'(X/S)\to
\operatorname{Pic}'(X\times_SS'/S')\mathrel{\substack{\longrightarrow\\[-0.6ex]\longrightarrow}}
\operatorname{Pic}'(X\times_SS''/S'')}
\]

\src{147}
(où $\symup{S''=S'\times_SS'}$) est exact.

Le premier énoncé résulte, grâce aux propriétés élémentaires de la descente
fidèlement plate, de la remarque générale suivante. Si $\symup{f:X\to S}$
est un morphisme tel que $\symup{\mathcal O_S\overset{\sim}{\to}f_*(\mathcal O_X)}$,
alors le foncteur $\symup{\mathfrak F\rightsquigarrow f^*(\mathfrak F)}$, de la catégorie des
Modules localement libres de type fini sur $\symup{S}$ dans la catégorie des Modules
localement libres de type fini sur $\symup{X}$, est pleinement fidèle, et l'image
essentielle est formée des Modules $\symup{\mathfrak G}$ sur $\symup{X}$ tels que
$\symup{f_*(\mathfrak G)}$ soit localement libre, et que l'homomorphisme canonique
\[
\symup{f^*\bigl(f_*(\mathfrak G)\bigr)\to\mathfrak G}
\]
soit un isomorphisme. Le deuxième énoncé a été prouvé par la théorie de la
descente dans [4], I, B, 4.

Les résultats de 2.2 s'énoncent aussi ainsi :

\result{COROLLAIRE}{2.3} Sous les conditions de 2.2, supposées vérifiées après tout changement de base, l'homomorphisme canonique
(1.9) $\symup{\mathop{\text{\uwave{Pic}}}\nolimits'_{X/S}\to\mathop{\text{\uwave{Pic}}}\nolimits_{X/S}}$ est
injectif, et même bijectif si $\symup{X}$ admet une section localement sur $\symup{S}$. (Donc
dans ce dernier cas, il y a identité entre le groupe de Picard relatif
$\symup{\operatorname{Pic}(X/S)}$ et le groupe de Picard relatif restreint
$\symup{\operatorname{Pic}'(X/S)}$.)

Conjuguant avec 2.1, on trouve donc :

\result{COROLLAIRE}{2.4} Sous les conditions de 2.2, on a une suite exacte
\[
\symup{0\to\operatorname{Pic}(S)\to\operatorname{Pic}(X)
\to\operatorname{Pic}(X/S).}
\]
Lorsque $\symup{X}$ admet une section sur $\symup{S}$, le dernier homomorphisme est surjectif,
i. e. on a alors un isomorphisme
\[
\symup{\operatorname{Pic}(X/S)\overset{\sim}{\to}
\operatorname{Pic}(X)/\operatorname{Pic}(S).}
\]

\result{\uline{Remarque}}{2.5} Soit $\symup{f:X\to S}$ un morphisme tel que
$\symup{\mathcal O_S\overset{\sim}{\to}f_*(\mathcal O_X)}$, et soit $\symup{g}$ une section de
$\symup{X}$ sur $\symup{S}$. Soit $\symup{\mathcal L}$ un Module inversible sur $\symup{X}$, appelons
$\symup{g}$-\uline{rigidification} de $\symup{\mathcal L}$ un isomorphisme
$\symup{\mathcal O_S\overset{\sim}{\to}g^*(\mathcal L)}$, et appelons \uline{Module
inversible} $\symup{g}$-\uline{rigidifié} un Module inversible $\symup{\mathcal L}$ sur $\symup{X}$ muni d'une
$\symup{g}$-rigidification. Tout automorphisme d'une telle structure est trivial, et
$\symup{\operatorname{Pic}'(X/S)}$ s'identifie à l'ensemble des classes, à un
isomorphisme près, de Modules inversibles $\symup{g}$-rigidifiés sur $\symup{X}$. (C'est ce fait qui a
permis d'utiliser la théorie de la descente pour prouver 2.2, (ii).) Cela donne
une nouvelle interprétation de $\symup{\operatorname{Pic}(X/S)}$, du moins lorsque $\symup{f}$
est de plus quasi-compact et séparé, donc
$\symup{\operatorname{Pic}(X/S)\overset{\sim}{\to}\operatorname{Pic}'(X/S)}$ par 2.3.

\src{148}
\result{\uline{Remarque}}{2.6} Soit $\symup{f:X\to S}$ un morphisme comme dans
2.2, et soit $\symup{S'\to S}$ un morphisme fidèlement plat et
quasi-compact tel qu'il existe un $\symup{S}$-morphisme $\symup{S'\to X}$, i. e. tel
qu'il existe une section de $\symup{X'=X\times_SS'}$ sur $\symup{S'}$. Soit
$\symup{S''=S'\times_SS'}$, $\symup{X''=X\times_SS''}$, et considérons la suite exacte
\[
\symup{\operatorname{Pic}(X/S)\to\operatorname{Pic}(X'/S')
\mathrel{\substack{\longrightarrow\\[-0.6ex]\longrightarrow}}\operatorname{Pic}(X''/S'').}
\]
Appliquant 2.4, on trouve une suite exacte.
\[
\symup{\operatorname{Pic}(X/S)\to
\operatorname{Pic}(X')/\operatorname{Pic}(S')
\mathrel{\substack{\longrightarrow\\[-0.6ex]\longrightarrow}}
\operatorname{Pic}(X'')/\operatorname{Pic}(S''),}
\]
en particulier, tout élément du groupe de Picard relatif ``provient'' déjà
d'un élément de $\symup{\operatorname{Pic}(X')}$. Cela donne donc une simplification
substantielle pour la description du groupe de Picard relatif donnée dans le
numéro précédent, et de même pour le foncteur de Picard de $\symup{X}$ sur $\symup{S}$ si les conditions de 2.2 sont vérifiées après tout changement de base, puisque,
pour tout $\symup{T}$ sur $\symup{S}$, on peut appliquer ce qui précède à
$\symup{X\times_ST/T}$ et au morphisme $\symup{T'=S'\times_ST\to T}$. Si par
exemple $\symup{f}$ lui-même est fidèlement plat, on peut prendre $\symup{S'=X}$, ce qui permet,
lorsque $\symup{f}$ est de plus de type fini (resp. simple, etc.) de se borner, dans la
description du foncteur de Picard relatif
$\symup{T\rightsquigarrow\operatorname{Pic}(X\times_ST/T)}$, à des changements de base
$\symup{T'\to T}$ qui sont de type fini (resp. simple, etc.). Lorsque $\symup{f}$
est projectif et plat, $\symup{S}$ localement noethérien, on prouve qu'on peut prendre
ci-dessus un $\symup{S'\to S}$ tel que $\symup{S'}$ soit somme directe de
\uline{revêtements plats} $\symup{S'_i}$ d'ouverts $\symup{S_i}$ de $\symup{S}$ couvrant $\symup{S}$ ; si $\symup{f}$ est même
séparable, on peut prendre $\symup{S'_i}$ étale sur $\symup{S_i}$.

\fsection{3}{\normalfont\uwave{Le théorème principal d'existence} : \uwave{énoncé}.}

On ne dispose pas, même à titre conjectural, d'un énoncé d'existence de
préschémas de Picard, englobant tous les cas connus. Une condition ``pratiquement
nécessaire'', si on peut dire, est que $\symup{f:X\to S}$ soit \uline{propre}
(assurant des propriétés de finitude essentielles) et \uline{plat}. Ces conditions ne
sont pas suffisantes, même si $\symup{S}$ est le spectre de l'algèbre des nombres duaux
$\symup{k[t]/(t^2)}$ sur un corps $\symup{k}$ (disons le corps $\symup{\text{\uwave{C}}}$ des nombres
complexes), et $\symup{X}$ de dimension 1. Au moment d'écrire le présent exposé, les
théorèmes d'existence les plus importants pour le préschéma de Picard sont
déduits du théorème suivant :

\result{THÉORÈME}{3.1} Soit $\symup{f:X\to S}$ un morphisme de préschémas
localement noethériens. On suppose

\noindent (i) $\symup{f}$ projectif

\src{149}
\noindent (ii) $\symup{f}$ plat

\noindent (iii) les fibres géométriques de $\symup{f}$ sont intègres.

Sous ces conditions, $\symup{\underline{\operatorname{Pic}}_{X/S}}$ existe.

La démonstration, qui sera esquissée dans les deux numéros qui suivent,
montrera en même temps ceci : Soit $\symup{\xi}$ la section de
$\symup{\underline{\operatorname{Pic}}_{X/S}}$ correspondant à un faisceau très ample
$\symup{\mathcal O_X(1)}$ sur $\symup{X/S}$ (i. e. induit par une immersion projective
$\symup{X\to P(\mathfrak E)}$) ; il existe alors une partie ouverte
$\symup{U}$ de $\symup{\underline{\operatorname{Pic}}_{X/S}}$, réunion disjointe de parties ouvertes
quasi-projectives sur $\symup{S}$, telle que $\symup{U}$ soit stable par la translation par
$\symup{\xi}$, et que $\symup{\underline{\operatorname{Pic}}_{X/S}}$ soit réunion croissante des ouverts
$\symup{U-n\xi}$ (tous isomorphes à $\symup{U}$). Il en résulte en particulier que sous les
conditions de 3.1, $\symup{\underline{\operatorname{Pic}}_{X/S}}$ est \uline{séparé} sur $\symup{S}$.

\result{\uline{Remarque}}{3.2} On voit sur des exemples (avec $\symup{S}$ le spectre d'un
anneau de valuation discrète, et $\symup{X}$ de dimension relative 1 sur $\symup{S}$, par
exemple), que si dans 3.1 on omet l'hypothèse (iii) en la remplaçant par
l'hypothèse plus faible que, pour tout $\symup{s\in S}$, l'homomorphisme
$\symup{k(s)\to H^0(X_s,\mathcal O_{X_s})}$ est un isomorphisme, alors $\symup{\underline{\operatorname{Pic}}_{X/S}}$ n'est pas nécessairement
séparé sur $\symup{S}$ ; tant dans le cas où les fibres géométriques de $\symup{f}$ sont
réduites, mais où une fibre géométrique générique intègre ``éclate'' par
spécialisation en deux composantes irréductibles, que dans le cas où les fibres
géométriques de $\symup{f}$ sont irréductibles, mais où une fibre géométrique générique
intègre se spécialise en une ``fibre multiple''. Le premier cas se présente par
exemple avec une conique dégénérant en deux droites concourantes, un exemple
du deuxième m'a été fourni par D. MUMFORD, avec une courbe elliptique dégénérant
en une courbe elliptique double. Ces exemples sont valables en toute
caractéristique.

\result{\uline{Remarque}}{3.3} Sous les conditions de 3.1, j'ignore si
$\symup{\underline{\operatorname{Pic}}_{X/S}}$ est réunion disjointe d'ouverts qui sont de type
fini, donc quasi-projectifs, sur $\symup{S}$. On notera que la considération des polynômes
de Hilbert $\symup{Q\in\text{\uwave{Q}}[t]}$ permet, comme dans le cas des schémas de Hilbert
([4], IV), de donner une décomposition de $\symup{\underline{\operatorname{Pic}}_{X/S}}$ en somme
disjointe d'ouverts $\symup{\underline{\operatorname{Pic}}^{Q}_{X/S}}$, et il semble plausible que
ces derniers sont quasi-projectifs sur $\symup{S}$ ; c'est ce qu'on verra du moins dans
le prochain exposé lorsque $\symup{f}$ est un morphisme simple. On fera attention que
si on remplace l'hypothèse (i) par l'hypothèse : $\symup{X}$ est projectif \uline{localement}
au-dessus de $\symup{S}$ (ce qui est suffisant pour la validité de 3.1, puisque la
question d'existence de $\symup{\underline{\operatorname{Pic}}_{X/S}}$ est visiblement locale sur
$\symup{S}$), il est facile par contre

\src{150}
de donner des exemples où $\symup{\underline{\operatorname{Pic}}_{X/S}}$ \uline{contient des composantes
connexes qui ne sont pas de type fini sur} $\symup{S}$. Soit par exemple $\symup{X_0}$ une
variété algébrique non singulière projective sur un corps $\symup{k}$ algébriquement
clos, munie d'un automorphisme $\symup{u}$ et d'un élément $\symup{\xi}$ dans le groupe de
Néron-Severi de $\symup{X_0}$, tels que les $\symup{u^n(\xi)}$ soient deux à deux distincts.
On peut par exemple prendre pour $\symup{X_0}$ le produit d'une courbe elliptique $\symup{E}$
par elle-même, pour $\symup{u}$ l'automorphisme
$\symup{(x,y)\rightsquigarrow(x,y+x)}$ de $\symup{E\times E}$. Soit $\symup{S}$ la réunion de deux courbes irréductibles non
singulières se coupant en deux points $\symup{a}$ et $\symup{b}$. Il y a sur $\symup{S}$ un revêtement
principal connexe $\symup{P}$ de groupe $\symup{\text{\uwave{Z}}}$, et utilisant les opérations de
$\symup{\text{\uwave{Z}}}$ sur $\symup{X_0}$ définies par $\symup{u}$, on en conclut un fibré associé sur $\symup{S}$,
de fibre $\symup{X_0}$ (trivial sur $\symup{S-a}$ et $\symup{S-b}$), en fait un schéma \uline{abélien} sur $\symup{S}$
dans le cas particulier envisagé. On voit facilement que
$\symup{\underline{\operatorname{Pic}}_{X/S}}$, qui est aussi le fibré associé à $\symup{P}$ et aux
opérations de $\symup{\text{\uwave{Z}}}$ sur $\symup{\underline{\operatorname{Pic}}_{X_0/k}}$ via $\symup{u}$, contient
une composante connexe isomorphe à
$\symup{P\times\underline{\operatorname{Pic}}^{0}_{X_0/k}}$ (où $\symup{\underline{\operatorname{Pic}}^{0}}$ dénote la composante connexe de l'élément neutre
dans $\symup{\underline{\operatorname{Pic}}}$), laquelle n'est pas de type fini sur $\symup{S}$. [Il se
produit également des phénomènes analogues dans divers cas de préschémas de
Picard non séparés sur $\symup{S}$, comme envisagés dans 3.2.]

\fsection{4}{\normalfont\uwave{Diviseurs de Cartier relatifs et fibrés projectifs}.}

Nous n'aurons à utiliser que des diviseurs positifs, et omettrons ce
qualificatif supplémentaire dans la suite du numéro.

Soit $\symup{X}$ un préschéma. Un \uline{diviseur de Cartier} sur $\symup{X}$, ou simplement
diviseur sur $\symup{X}$ pour simplifier, est un sous-préschéma fermé $\symup{D}$ de $\symup{X}$, défini
par un Idéal $\symup{\mathfrak J}$ qui est un Module \uline{inversible}, i. e. engendré localement
par une section \uline{non diviseur de zéro} de $\symup{\mathcal O_X}$. À $\symup{D}$ nous associons un
Module inversible, à savoir
\[
\symup{\mathcal L(D)=\mathfrak J^{-1},}
\]
et l'injection canonique $\symup{\mathfrak J\to\mathcal O_X}$ donne un
homomorphisme canonique
\[
\symup{s_D:\mathcal O_X\to\mathfrak J^{-1}=\mathcal L(D),
\qquad\text{i. e.}\quad s_D\in\Gamma(X,\mathcal L(D)).}
\]
D'ailleurs, la donnée d'un diviseur est essentiellement \uline{équivalente} à la
donnée d'un Module inversible $\symup{\mathcal L}$ sur $\symup{X}$, muni d'une section $\symup{s}$ qui
soit partout non diviseur de zéro, en associant inversement à un tel couple
$\symup{(\mathcal L,s)}$ le ``diviseur'' de $\symup{s}$, noté $\symup{\operatorname{div}(s)}$. Pour un
$\symup{\mathcal L}$ inversible donné sur $\symup{X}$, l'ensemble des

\src{151}
diviseurs $\symup{D}$ définissant $\symup{\mathcal L}$ est en correspondance biunivoque avec
l'ensemble quotient $\symup{\Gamma(X,\mathcal L)^*/\Gamma(X,\mathcal O_X^*)}$, où $\symup{\Gamma(X,\mathcal L)^*}$ désigne la partie de $\symup{\Gamma(X,\mathcal L)}$ formée
des sections qui sont partout non diviseur de zéro.

Supposons maintenant que l'on ait un morphisme localement de type fini
$\symup{f:X\to S}$, et supposons pour simplifier $\symup{S}$ localement noethérien.
Soient $\symup{\mathfrak J}$ un Idéal cohérent sur $\symup{X}$, $\symup{D}$ le sous-schéma de $\symup{X}$ qu'il
définit, $\symup{x\in X}$, et $\symup{s=f(x)}$. On montre que les conditions suivantes sont
équivalentes :

\noindent (i) $\symup{\mathfrak J}$ est inversible en $\symup{x}$ (i. e. $\symup{\mathfrak J_x}$ engendré
par un élément régulier de $\symup{\mathcal O_{X,x}}$) et $\symup{X}$ et $\symup{D}$ sont plats sur $\symup{S}$ en $\symup{x}$.

\noindent (ii) $\symup{X}$ et $\symup{D}$ sont plats sur $\symup{S}$ en $\symup{x}$, et $\symup{D_s}$ est un diviseur
de Cartier sur la fibre $\symup{X_s}$ en le point $\symup{x}$.

\noindent (iii) $\symup{X}$ est plat sur $\symup{S}$ en $\symup{x}$, et $\symup{\mathfrak J_x}$ est engendré par
un élément $\symup{f_x}$ induisant sur $\symup{X_s}$ un germe non diviseur de zéro.

On dit alors que $\symup{D}$ est un \uline{diviseur de Cartier relatif} sur $\symup{X/S}$, (ou
simplement diviseur relatif sur $\symup{X/S}$), au point envisagé. On notera sur (i)
qu'alors $\symup{D}$ est également un diviseur relatif aux points voisins de $\symup{x}$, donc
si $\symup{X}$ et $\symup{D}$ sont plats sur $\symup{S}$, et $\symup{D}$ propre sur $\symup{S}$, alors l'ensemble des
$\symup{s\in S}$ tels que $\symup{D_s}$ soit un diviseur de Cartier dans $\symup{X_s}$ (i. e. tels que
$\symup{D}$ soit un diviseur de Cartier relatif aux points de $\symup{X_s}$) est une partie
ouverte de $\symup{S}$. D'autre part, on a fait ce qu'il fallait dans la définition
précédente pour que la notion de diviseur de Cartier relatif soit stable par
un changement de base quelconque $\symup{S'\to S}$. Considérons alors,
l'ensemble $\symup{\operatorname{Div}(X/S)}$ des diviseurs relatifs sur $\symup{X/S}$, puis le
foncteur contravariant en $\symup{S'}$ variable sur $\symup{S}$ défini par
\[
\symup{\text{\uwave{Div}}_{X/S}(S')=
\operatorname{Div}(X\times_SS'/S').}
\]

Supposons $\symup{X}$ plat et propre sur $\symup{S}$. Alors par la caractérisation (ii) des
diviseurs de Cartier relatifs, $\symup{\text{\uwave{Div}}_{X/S}}$ peut être considéré
comme un sous-foncteur du foncteur $\symup{\text{\uwave{Hilb}}_{X/S}}$ défini dans
[4], IV, et le morphisme d'inclusion
\[
\symup{\text{\uwave{Div}}_{X/S}\to\operatorname{Hilb}_{X/S}}
\]
est ``représentable par des immersions ouvertes'', (cf. [5], IV, 3.13) en vertu
des remarques qui précèdent. Utilisant le théorème d'existence principal de
[4], IV, on trouve :

\src{152}
\result{PROPOSITION}{4.1} Supposons $\symup{f:X\to S}$ projectif et plat.
Alors le foncteur $\symup{\text{\uwave{Div}}_{X/S}}$ est représentable, et, de façon
précise, est représentable par un ouvert de $\symup{\underline{\operatorname{Hilb}}_{X/S}}$.

Utilisant (pour un faisceau $\symup{\mathcal O_X(1)}$ très ample donné sur $\symup{X/S}$) la
décomposition canonique de $\symup{\underline{\operatorname{Hilb}}_{X/S}}$ en somme des ouverts
$\symup{\underline{\operatorname{Hilb}}^{Q}_{X/S}}$ correspondants aux polynômes de Hilbert
$\symup{Q\in\text{\uwave{Q}}[t]}$, on en conclut d'ailleurs une décomposition analogue
\[
\symup{\underline{\operatorname{Div}}_{X/S}=
\coprod\limits_{Q\in\text{\uwave{Q}}[t]}\underline{\operatorname{Div}}^{Q}_{X/S}}
\]
en somme d'ouverts disjoints \uline{quasi-projectifs} sur $\symup{S}$.

Utilisant l'application $\symup{D\rightsquigarrow\mathcal L(D)}$, on trouve d'autre part un
homomorphisme fonctoriel
\begin{equation}
\symup{\text{\uwave{Div}}_{X/S}\to\text{\uwave{Pic}}_{X/S}
\tag{+}}
\end{equation}
que nous nous proposons d'étudier ; il apparaîtra qu'il est relativement
représentable ([5], IV, 3) sous des conditions assez générales. Partons donc
avec un élément $\symup{\xi}$ de $\symup{\text{\uwave{Pic}}_{X/S}(S')}$, en supposant pour
simplifier les notations que $\symup{S'=S}$, et montrons que le sous-foncteur
correspondant de $\symup{\text{\uwave{Div}}_{X/S}}$ est représentable. Prenons d'abord
le cas où $\symup{\xi}$ est défini par un Module inversible $\symup{\mathcal L}$ sur $\symup{X}$. Nous
supposerons $\symup{X}$ propre et plat sur $\symup{S}$, et que les fibres géométriques de $\symup{X}$
sur $\symup{S}$ sont intègres, ce qui implique aussi ([2], III, par. 7) que l'on a
$\symup{\mathcal O_S\overset{\sim}{\to}f_*(\mathcal O_X)}$ et que cette relation reste
valable après tout changement de base $\symup{S'\to S}$. Alors les
diviseurs de Cartier relatifs $\symup{D}$ sur $\symup{X/S}$ tels que $\symup{\mathcal L(D)}$ et
$\symup{\mathcal L}$ définissent le même élément de
$\symup{\operatorname{Pic}(X/S)=\text{\uwave{Pic}}_{X/S}(S)}$, i. e. en vertu de 2.4,
tels que $\symup{\mathcal L(D)}$ et $\symup{\mathcal L}$ soient isomorphes localement au-dessus
de $\symup{S}$, sont en correspondance biunivoque avec les sections du faisceau
quotient $\symup{f_*(\mathcal L)^*/\mathcal O_S^*}$. Cette correspondance est compatible
avec les changements de base. D'autre part, des considérations générales du
type ``Künneth'' de loco citato (cf. aussi [6]) montrent que la propreté de
$\symup{X/S}$ et la platitude de $\symup{\mathcal L}$ sur $\symup{S}$ impliquent l'existence d'un Module
cohérent $\symup{\mathfrak Q}$ sur $\symup{S}$, défini à un isomorphisme unique près, et un
isomorphisme de faisceaux
\[
\symup{f_*(\mathcal L)\overset{\sim}{\to}
\underline{\operatorname{Hom}}_{\mathcal O_S}(\mathfrak Q,\mathcal O_S),}
\]
la formation de $\symup{\mathfrak Q}$ étant par ailleurs
compatible avec les changements de base. Ici $\symup{f_*(\mathcal L)^*}$ désigne le
sous-faisceau d'ensembles de $\symup{f_*(\mathcal L)}$ dont les sections sur $\symup{U}$ sont
les sections de $\symup{\mathcal L}$ sur $\symup{f^{-1}(U)}$ qui définissent des diviseurs de

\src{153}
Cartier relatifs sur $\symup{f^{-1}(U)/U}$, i. e. qui induisent des sections non
diviseurs de $\symup{0}$ sur les $\symup{X_s}$ ($\symup{s\in U}$). Utilisant l'hypothèse que les
fibres $\symup{X_s}$ sont intègres, cela signifie simplement que les sections induites
sur les fibres $\symup{X_s}$ ne sont pas identiquement nulles, ou encore en termes
d'homomorphismes locaux $\symup{\mathfrak Q\to\mathcal O_S}$, que lesdits
homomorphismes soient surjectifs (NAKAYAMA). Cela montre que l'ensemble des
sections de $\symup{f_*(\mathcal L)^*/\mathcal O_S^*}$ est en correspondance
biunivoque avec l'ensemble des \uline{Modules quotients inversibles de} $\symup{\mathfrak Q}$,
ou encore, par définition du fibré projectif $\symup{P(\mathfrak Q)}$ associé au
Module cohérent $\symup{\mathfrak Q}$ (cf. [5], V, 2), avec l'ensemble des sections de
$\symup{P(\mathfrak Q)}$ sur $\symup{S}$. Cette description est compatible avec la
formation des images inverses, et on obtient par suite le théorème ci-après.

\result{THÉORÈME}{4.3} Soient $\symup{f:X\to S}$ un morphisme propre et
plat à fibres géométriques intègres, avec $\symup{S}$ localement noethérien,
$\symup{\mathcal L}$ un Module inversible sur $\symup{X}$. Pour tout $\symup{S'}$ sur $\symup{S}$, soit
$\symup{T(S')}$ l'ensemble des diviseurs relatifs $\symup{D}$ sur $\symup{X\times_SS'/S'}$ tels que
$\symup{\mathcal L(D)}$ soit isomorphe à
$\symup{\mathcal L\otimes_{\mathcal O_S}\mathcal O_{S'}}$ localement au-dessus de $\symup{S'}$
(i. e. tels que $\symup{\mathcal L(D)}$ et
$\symup{\mathcal L\otimes_{\mathcal O_S}\mathcal O_{S'}}$ définissent le même élément
de $\symup{\operatorname{Pic}(X\times_SS'/S')}$). Alors il existe un Module cohérent
$\symup{\mathfrak Q}$ sur $\symup{S}$, déterminé à un isomorphisme unique près, tel que le
foncteur $\symup{T}$ soit représentable par le fibré projectif $\symup{P(\mathfrak Q)}$.

\result{COROLLAIRE}{4.4} Si on suppose même $\symup{f}$ projectif, alors
l'homomorphisme fonctoriel
$\symup{\text{\uwave{Div}}_{X/S}\to\text{\uwave{Pic}}_{X/S}}$ est
représentable par morphismes projectifs.

En fait, lorsque $\symup{X}$ admet une section (resp. localement une section) sur $\symup{S}$,
alors l'homomorphisme précédent est représentable par des fibrés projectifs
(resp. par des fibrés projectifs locaux), grâce à 4.3 et 2.1. Dans le cas où
$\symup{f}$ est quasi-projectif, on se ramène facilement au cas précédent par une
méthode de descente, utilisant des quasi-sections finies et plates de $\symup{X}$
localement sur $\symup{S}$.

\result{\uline{Remarque}}{4.5} Sous les conditions de 4.3, le Module $\symup{\mathfrak Q}$ sur
$\symup{S}$ n'est en général pas localement libre, ce qui s'exprime par le fait que la
dimension des fibres réduites de $\symup{\mathfrak Q}$, i. e. celle des
$\symup{H^0(X_s,\mathcal L_s)}$ pour $\symup{s\in S}$ variable, peut faire des \uline{sauts}. Étant
donné un Module cohérent $\symup{\mathfrak Q}$ sur le préschéma localement noethérien
$\symup{S}$, on vérifie facilement que pour un $\symup{s\in S}$ donné, $\symup{\mathfrak Q}$ est libre
en $\symup{s}$ si et seulement si $\symup{P(\mathfrak Q)}$ est plat sur $\symup{S}$ en les
points au-dessus de $\symup{s}$, (auquel cas il sera donc même simple sur $\symup{S}$ en les
points au-dessus de $\symup{s}$) ; en l'occurrence, $\symup{\mathfrak Q}$ étant défini en termes
de $\symup{\mathcal L}$ comme ci-dessus, cela signifie aussi que la formation de
l'image directe $\symup{f_*(\mathcal L)}$ ``\uline{commute avec le changement de base au}

\src{154}
\uline{voisinage de} $\symup{s}$'', ou encore que
$\symup{f_*(\mathcal L)_s\to H^0(X_s,\mathcal L_s)}$ est surjectif. Il en
sera ainsi par exemple si $\symup{H^1(X_s,\mathcal L_s)=0}$. Sous réserve d'existence
des préschémas en jeu, ces critères s'appliquent en particulier à la situation
universelle $\symup{\underline{\operatorname{Div}}_{X/S}\to\underline{\operatorname{Pic}}_{X/S}}$,
et donnent une condition nécessaire et suffisante, resp. suffisante, pour que
ce morphisme soit simple en un point donné de $\symup{\underline{\operatorname{Div}}_{X/S}}$.

\fsection{5}{\normalfont\uwave{Démonstration du théorème principal d'existence}.}

Sous les conditions de 3.1, choisissons un Module $\symup{\mathcal O_X(1)}$ très ample
sur $\symup{X/S}$, soit $\symup{\xi}$ l'élément correspondant de
$\symup{\operatorname{Pic}(X/S)}$. Posons pour abréger
$\symup{\text{\uwave{P}}(S')=\operatorname{Pic}(X\times_SS'/S')}$, et supposons d'abord pour simplifier que $\symup{X/S}$ admette une section. Soit
$\symup{\text{\uwave{P}}^+(S')}$ la partie de $\symup{\text{\uwave{P}}(S')}$ formée des classes des $\symup{\mathcal L}$ inversibles
sur $\symup{X\times_SS'}$ tels que
\[
\begin{gathered}
\symup{R^if'_*(\mathcal L(n))=0}\quad\text{pour $\symup{i>0}$ et tout $\symup{n\geq0}$}\\
\symup{f'_*(\mathcal L(n))\neq0}\quad\text{en tout point de $\symup{S'}$ et pour tout $\symup{n\geq0}$}.
\end{gathered}
\]
Ce sont là des conditions stables par changement de base, donc qui définissent
un sous-foncteur $\symup{\text{\uwave{P}}^+}$ de $\symup{\text{\uwave{P}}}$, évidemment stable par translation par $\symup{\xi}$.
Utilisant les ``théorèmes A et B'' de SERRE ([2], III, 2) et les généralités
([5], IV, 5), on voit facilement que $\symup{\text{\uwave{P}}}$ est représentable si et seulement si
$\symup{\text{\uwave{P}}^+}$ l'est, et alors $\symup{\text{\uwave{P}}^+}$ sera représentable par un ouvert $\symup{U}$ du préschéma
$\symup{\underline{\operatorname{Pic}}_{X/S}}$ représentant $\symup{\text{\uwave{P}}}$, et ce dernier sera réunion
croissante des ouverts $\symup{U-n\xi}$.

Soit pour abréger $\symup{\text{\uwave{D}}=\text{\uwave{Div}}_{X/S}}$, et soit $\symup{\text{\uwave{D}}^+}$ l'image inverse
de $\symup{\text{\uwave{P}}^+}$ par le morphisme canonique $\symup{\text{\uwave{D}}\to\text{\uwave{P}}}$. Donc on a un
morphisme
\begin{equation}
\symup{\text{\uwave{D}}^+\to\text{\uwave{P}}^+,
\tag{+}}
\end{equation}
et on sait déjà que $\symup{\text{\uwave{D}}^+}$ est représentable par un ouvert $\symup{D^+}$ du préschéma
$\symup{D=\underline{\operatorname{Div}}_{X/S}}$ (dont l'existence est garantie par 4.1). Il
résulte facilement de 4.3 que le morphisme (+) est représentable par morphismes
\uline{simples projectifs surjectifs} (et de façon précise par des \uline{fibrés projectifs
associés à des Modules localement libres} partout non nuls) : cela tient au fait
que, lorsque $\symup{\mathcal L}$ sur $\symup{X\times_SS'}$ est, comme au début du présent
numéro, alors $\symup{f'_*(\mathcal L)}$ est un Module \uline{localement libre} $\symup{\neq0}$, dont la
formation commute au changement de base ; avec les notations de 4.3,
$\symup{\mathfrak Q}$ sera alors isomorphe au dual de $\symup{f'_*(\mathcal L)}$. Utilisant le
critère général ([5], IV, 4.7), nous allons en conclure la représentabilité de
$\symup{P^+}$. Dans loco citato, nous prendrons pour $\symup{\mathfrak C}$ l'ensemble des
morphismes fidèlement

\src{155}
plats et quasi-compacts de préschémas (qui sont bien des épimorphismes
effectifs en vertu de [4], I B). La condition (a) de loco citato, à savoir que
(+) est représentable par morphismes éléments de $\symup{\mathfrak C}$, est vérifiée
comme on vient de le voir ; de même la condition (b), qui signifie que le
foncteur $\symup{\text{\uwave{P}}^+}$ est compatible avec la descente fidèlement plate
quasi-compacte, ce qui est immédiat. Reste à prouver la condition (c) de loco
citato, à savoir que l'équivalence $\symup{R}$ dans le préschéma $\symup{D^+}$ déduite du
morphisme $\symup{\mathfrak C}$-représentable (+) est $\symup{\mathfrak C}$-effective, i. e. est
effective et telle que $\symup{D^+\to D^+/R}$ soit $\symup{\in\mathfrak C}$.
Pour ceci on note d'abord que les ouverts $\symup{D^{+Q}}$ de $\symup{D^+}$ correspondants aux
divers polynômes de Hilbert virtuels $\symup{Q\in\text{\uwave{Q}}[t]}$ sont stables par $\symup{R}$
(car les fibres de $\symup{R}$ sont connexes), ce qui nous ramène à montrer que pour
tout $\symup{Q}$, la relation d'équivalence $\symup{R^Q}$ induite dans $\symup{D^{+Q}}$ est
$\symup{\text{\uwave{S}}}$-effective. Or maintenant $\symup{D^{+Q}}$ est quasi-projectif, et d'autre
part la relation d'équivalence $\symup{R^Q}$ est projective et plate. On est donc sous
les conditions d'application de [4], III, 6.1, qui implique le résultat voulu.

Dans le cas général où $\symup{X/S}$ n'admet pas nécessairement de section, on se
ramène facilement au cas précédent par la technique de descente, ou on reprend
la démonstration précédente avec la modification qui s'impose dans la
définition de $\symup{\text{\uwave{P}}^+}$.

\result{\uline{Remarques}}{5.1} La méthode suivie est essentiellement celle de MATSUSAKA
pour la construction projective des variétés de Picard. Le résultat invoqué de
[4], III pour la possibilité de passage au quotient effectif peut se déduire
facilement aussi du théorème d'existence des schémas de Hilbert, (cf. par
exemple [6]) (Classiquement, ces quotients se construisaient par utilisation
des coordonnées de Chow). On notera que la formation de l'ouvert
$\symup{\underline{\operatorname{Pic}}^{+}_{X/S}}$ de $\symup{\underline{\operatorname{Pic}}_{X/S}}$ et sa
décomposition en ouverts $\symup{\underline{\operatorname{Pic}}^{+Q}_{X/S}}$ quasi-projectifs sur
$\symup{S}$, suivant les polynômes de Hilbert pour les \uline{diviseurs} définissant les
Modules inversibles envisagés, est compatible avec le changement de base (ce
qui permet l'application de la technique de descente).

\result{\uline{Remarques}}{5.2} Il n'est pas exclu que $\symup{\underline{\operatorname{Pic}}_{X/S}}$ existe
chaque fois que $\symup{f:X\to S}$ est propre et plat et tel que les
homomorphismes $\symup{k(s)\to H^0(X_s,\mathcal O_{X_s})}$
$\symup{(s\in S)}$ soient des isomorphismes (cette dernière condition signifiant alors aussi que
$\symup{\mathcal O_S\overset{\sim}{\to}f_*(\mathcal O_X)}$ et que cette relation reste
valable après tout changement de base $\symup{S'\to S}$). C'est du moins
ce qu'on prouve dans le cadre des espaces analytiques lorsque $\symup{f}$ est de plus
projectif, par une méthode différente (celle de
CHOW sauf erreur) expliquée
dans [5], IX, 3.1. Dans cette
méthode, le passage au

\src{156}
quotient par une relation d'équivalence \uline{propre} et plate dans un ouvert du
schéma des diviseurs, est remplacé par un passage au quotient par le groupe
projectif dans le schéma des immersions de $\symup{X}$ dans $\symup{\text{\uwave{P}}^r_S}$. Cette
méthode pourra probablement s'adapter au cas des schémas, en utilisant les
résultats de MUMFORD sur le passage au quotient par le groupe projectif [8] ;
pour l'instant, il n'y a de démonstration écrite que lorsque $\symup{X}$ a ``beaucoup
de sections locales'' au-dessus de $\symup{S}$, par exemple lorsque $\symup{X}$ est séparable
sur un anneau local complet. En principe, la méthode en question serait de
portée plus générale, puisqu'elle donne également l'existence de préschémas de
Picard dans des cas où ceux-ci ne sont pas séparés, et où la première méthode
échoue donc nécessairement. (Techniquement, la difficulté provient du fait que
lorsque les fibres géométriques de $\symup{f}$ ne sont plus intègres, alors le foncteur
envisagé dans 4.3 n'est plus représentable par le fibré projectif
$\symup{P(\mathfrak Q)}$ lui-même, mais par un \uline{ouvert} dudit, ce qui conduit à la
question délicate du passage au quotient par une relation d'équivalence plate
mais non propre.)

\result{\uline{Remarque}}{5.3} On notera que la démonstration donnée ici ne fait appel,
ni à la construction préalable des jacobiennes des courbes ou familles de
courbes, ni à la théorie des variétés abéliennes ou schémas abéliens, et par là
se distingue de façon essentielle d'exposés traditionnels comme dans le livre
de LANG [7] ou l'exposé de CHEVALLEY [1], qui suivent la voie esquissée par
A. WEIL. Même dans le cas des jacobiennes des courbes non singulières sur un
corps algébriquement clos (le corps des complexes disons), la construction
donnée ici pour la jacobienne est la seule connue qui en fournisse les
propriétés très fortes que nous avons prises comme définition au numéro 1,
(essentiellement celles de CHEVALLEY, mais en tenant compte de ``variétés de
paramètres'' à éléments nilpotents). Que la construction des schémas de Picard
doive précéder et non suivre la théorie des variétés abéliennes, est clair a
priori par le fait qu'en général les schémas de Picard ne sont ni ne se
ramènent à des schémas abéliens, comme on voit déjà dans le cas de courbes
singulières sur un corps algébriquement clos, où on trouve les ``jacobiennes
généralisées'' de Rosenlicht, qui ne sont pas des variétés abéliennes. De plus,
la théorie des variétés abéliennes, et plus généralement des schémas abéliens,
gagne beaucoup en simplicité une fois qu'on dispose d'une théorie des schémas
de Picard en général. En particulier, la théorie de dualité pour les schémas
abéliens, et notamment les résultats du type Cartier, deviennent à peu près
formels à partir de là (cf. par exemple [10]).

\src{157}
\result{\uline{Remarques}}{5.4} Le ``principe de compatibilité'' de Igusa pour la
jacobienne d'une courbe dégénérant en une courbe singulière, ne peut être bien
compris que comme un théorème d'existence du schéma de Picard d'un schéma
relatif en courbes $\symup{X/S}$, non nécessairement simple sur $\symup{S}$. C'est donc un cas
particulier du théorème d'existence principal 3.1 lorsque la courbe spécialisée
est intègre (i. e. en termes classiques, irréductible de multiplicité 1). On
notera que pour l'instant, le cas d'une courbe spéciale réductible (même
lorsque les composantes sont de multiplicité 1, i. e. lorsque la courbe spéciale
est séparable sur le corps résiduel) échappe aux théorèmes d'existence connus,
sauf dans le cas où on est sur un anneau de valuation discrète \uline{complet} à corps
résiduel algébriquement clos, comparer 5.2. Cette question d'existence se posera
certainement à l'occasion d'une construction géométrique, en théorie des
schémas, des ``compactifications'' de Baily-Satake des schémas modulaires des
courbes de genre $\symup{g}$. (Cette compactification est connue pour l'instant
seulement pour $\symup{g=1}$, grâce aux travaux de IGUSA).

\fsection{6}{\normalfont\uwave{Théorèmes d'existence relatifs}.}

Nous allons esquisser ici quelques cas utiles où l'existence de certains
schémas de Picard implique l'existence de certains autres, ce qui permet de
déduire du théorème principal 3.1 divers autres théorèmes d'existence.

\result{PROPOSITION}{6.1} Soit $\symup{f:X\to S}$ un morphisme plat
projectif tel que dans la factorisation de Stein $\symup{f=f''f'}$, le morphisme
$\symup{f':X\to S'}$ soit \uline{plat} et à fibres géométriques intègres (donc
satisfait aux hypothèses de 3.1), et le morphisme fini
$\symup{f'':S'\to S}$ soit plat. Alors $\symup{\underline{\operatorname{Pic}}_{X/S}}$ existe et
(avec les notations introduites dans [4], II, C 2) on a un isomorphisme
canonique
\[
\symup{\underline{\operatorname{Pic}}_{X/S}\overset{\sim}{\to}
\prod\limits_{S'/S}\underline{\operatorname{Pic}}_{X/S'}.}
\]

Pour le montrer, on établit d'abord un isomorphisme de foncteurs
\[
\symup{\text{\uwave{Pic}}_{X/S}\overset{\sim}{\to}
\prod\limits_{S'/S}\text{\uwave{Pic}}_{X/S'},}
\]
on utilise 3.1, qui implique que $\symup{\text{\uwave{Pic}}_{X/S'}}$ est représentable,
et on utilise la structure signalée dans le n\textsuperscript{o} 3 de
$\symup{\underline{\operatorname{Pic}}_{X/S'}}$ (impliquant que toute partie finie d'une fibre de
$\symup{\underline{\operatorname{Pic}}_{X/S'}}$ sur $\symup{S}$ est contenue dans un ouvert affine) pour
l'existence de $\symup{\prod\limits_{S'/S}\underline{\operatorname{Pic}}_{X/S'}}$.

\src{158}
Par exemple, lorsque $\symup{X}$ est un schéma \uline{somme} de schémas $\symup{X_i}$ sur $\symup{S}$
satisfaisant les conditions de 3.1, l'énoncé 6.1 se réduit à l'énoncé trivial
\[
\symup{\underline{\operatorname{Pic}}_{X/S}\overset{\sim}{\to}
\prod\limits_i\underline{\operatorname{Pic}}_{X_i/S}.}
\]

\result{COROLLAIRE}{6.2} Soit $\symup{f:X\to S}$ un morphisme projectif et
plat à fibres géométriques localement intègres (par exemple un morphisme
projectif et normal), alors $\symup{\underline{\operatorname{Pic}}_{X/S}}$ existe.

Dans ce cas d'ailleurs, $\symup{S'}$ est un revêtement étale de $\symup{S}$ (vrai chaque fois
que $\symup{f}$ est séparable, i. e. plat à fibres géométriques réduites), et on vérifie
que le théorème de structure énoncé dans le n\textsuperscript{o} 3 pour
$\symup{\underline{\operatorname{Pic}}_{X/S}}$ est encore valable, grâce à la structure analogue
de $\symup{\underline{\operatorname{Pic}}_{X/S'}}$.

D'autre part, un procédé de descente donne un théorème d'existence relatif,
dont la portée dépend d'ailleurs de la solution des questions de descente non
plate soulevées dans [4], I, A (c), et dont nous nous contentons d'expliciter
ici un cas particulier :

\result{PROPOSITION}{6.3} Soit $\symup{f:X\to S}$ un morphisme propre,
soient $\symup{X_1}$ et $\symup{X_2}$ deux sous-préschémas de $\symup{X}$ plats sur $\symup{S}$, définis par
deux Idéaux cohérents $\symup{\mathfrak J_1}$ et $\symup{\mathfrak J_2}$ tels que
$\symup{\mathfrak J_1\cap\mathfrak J_2=(0)}$ et que
$\symup{\mathcal O_X/(\mathfrak J_1+\mathfrak J_2)}$ soit plat sur $\symup{S}$ (i. e. le
sous-préschéma de $\symup{X}$ sup de $\symup{X_1}$ et $\symup{X_2}$ est $\symup{X}$, tandis que leur inf $\symup{Z}$
est plat sur $\symup{S}$). Supposons de plus que, pour tout $\symup{s\in S}$, les Hom $\symup{k(s)\to H^0(X_{i_s},\mathcal O_{X_{i_s}})}$,
$\symup{i=1,2}$, soient bijectifs. Alors l'homomorphisme naturel de foncteurs
\[
\symup{\text{\uwave{Pic}}_{X/S}\to
\text{\uwave{Pic}}_{X_1/S}\times\text{\uwave{Pic}}_{X_2/S}}
\]
est représentable par des morphismes affines, donc si
$\symup{\underline{\operatorname{Pic}}_{X_1/S}}$ et $\symup{\underline{\operatorname{Pic}}_{X_2/S}}$ existent, il en
est de même de $\symup{\underline{\operatorname{Pic}}_{X/S}}$, et le morphisme canonique
\[
\symup{\underline{\operatorname{Pic}}_{X/S}\to
\underline{\operatorname{Pic}}_{X_1/S}\times\underline{\operatorname{Pic}}_{X_2/S}}
\]
est affine.

Par descente fidèlement plate, on est ramené au cas où $\symup{Z}$ admet une section
sur $\symup{S}$, définissant donc des sections de $\symup{X}$, $\symup{X_1}$, $\symup{X_2}$ sur $\symup{S}$,
permettant d'éliminer les automorphismes dans les structures considérées comme
expliqué dans 2.5. La démonstration consiste alors à noter que la donnée d'un
Module inversible

\src{159}
``rigidifié'' $\symup{\mathcal L}$ sur $\symup{X}$ équivaut à la donnée d'un triple
$\symup{(\mathcal L_1,\mathcal L_2,u)}$, où $\symup{\mathcal L_i}$ est un Module inversible
``rigidifié'' sur $\symup{X_i}$, et $\symup{u}$ un isomorphisme de
$\symup{\mathcal L_1|Z}$ avec $\symup{\mathcal L_2|Z}$, compatible avec les rigidifications.
Il y a seulement à vérifier que pour $\symup{\mathcal L_1}$, $\symup{\mathcal L_2}$ fixés, la
donnée de $\symup{u}$ s'exprime par une section d'un schéma convenable sur $\symup{S}$ affine
sur $\symup{S}$, ce qui est facile. On conclut facilement de 6.3 :

\result{COROLLAIRE}{6.4} Soit $\symup{X}$ un schéma propre et séparable sur un corps
$\symup{k}$, et soient $\symup{X_i}$ les composantes irréductibles de $\symup{X}$. Si les
$\symup{\underline{\operatorname{Pic}}_{X_i/k}}$ existent, il en est de même de
$\symup{\underline{\operatorname{Pic}}_{X/k}}$, et le morphisme canonique
\[
\symup{\underline{\operatorname{Pic}}_{X/k}\to\prod\limits_i\underline{\operatorname{Pic}}_{X_i/k}}
\]
est affine.

Combiné à 6.2, cela montre par exemple l'existence de
$\symup{\underline{\operatorname{Pic}}_{X/k}}$ chaque fois que $\symup{X}$ est un schéma projectif et
séparable sur un corps $\symup{k}$. Lorsque $\symup{X}$ n'est plus séparable sur $\symup{k}$, on a
également un résultat de réduction, en utilisant un raisonnement dû à OORT
[11]. La méthode s'applique également sur un schéma de base quelconque (cas
utile par exemple pour démontrer dans l'exposé suivant le résultat de finitude
annoncé dans 3.3). Pour éviter un énoncé trop technique, bornons-nous au cas où
l'on est sur un corps de base :

\result{PROPOSITION}{6.5} Soient $\symup{X}$ un schéma propre sur un corps $\symup{k}$, $\symup{X_0}$
un sous-schéma ayant même ensemble sous-jacent (donc défini par un Idéal
nilpotent sur $\symup{X}$). Alors le morphisme fonctoriel
$\symup{\text{\uwave{Pic}}_{X/k}\to\text{\uwave{Pic}}_{X_0/k}}$ est représentable par morphismes affines. En particulier, si
$\symup{\underline{\operatorname{Pic}}_{X_0/k}}$ existe, il en est de même de
$\symup{\underline{\operatorname{Pic}}_{X/k}}$, et le morphisme
$\symup{\underline{\operatorname{Pic}}_{X/k}\to\underline{\operatorname{Pic}}_{X_0/k}}$ est
affine.

Combinant avec 6.4, on conclut facilement :

\result{COROLLAIRE}{6.6} Soit $\symup{X}$ un préschéma projectif sur un corps $\symup{k}$,
alors $\symup{\underline{\operatorname{Pic}}_{X/k}}$ existe.

\result{\uline{Remarque}}{6.6} Il est extrêmement plausible que pour tout schéma $\symup{X}$
\uline{propre} sur un corps $\symup{k}$, $\symup{\underline{\operatorname{Pic}}_{X/k}}$ existe.
Les résultats qui précèdent nous ramènent, dans cette question, au cas où $\symup{k}$
est algébriquement clos, et où $\symup{X}$ est intègre. On sait alors qu'il existe un
schéma intègre $\symup{X'}$ \uline{projectif} sur $\symup{k}$ et un morphisme dominant
$\symup{g:X'\to X}$ (lemme de Chow). Il suffirait donc de montrer que le

\src{160}
morphisme fonctoriel correspondant
$\symup{\text{\uwave{Pic}}_{X/k}\to\text{\uwave{Pic}}_{X'/k}}$ est représentable (et sans doute, représentable par morphismes affines),
puisqu'on sait par ailleurs que $\symup{\text{\uwave{Pic}}_{X'/k}}$ est représentable.
Cela soulève des questions de descente non plate non résolues pour l'instant.
On notera que lorsqu'on se borne à considérer la restriction du foncteur
$\symup{\text{\uwave{Pic}}_{X/k}}$ aux préschémas réduits, ($\symup{X}$ étant propre intègre
sur $\symup{k}$ algébriquement clos), on obtient bien un foncteur représentable, comme
l'a montré CHEVALLEY [1] dans le cas où $\symup{X}$ est normal, et SESHADRI [12] par une
méthode de descente dans le cas général. Mais avec nos notations, le schéma
construit par ces auteurs n'est pas $\symup{\underline{\operatorname{Pic}}_{X/k}}$, mais
$\symup{(\underline{\operatorname{Pic}}_{X/k})_{\mathrm{réd}}}$, le schéma réduit correspondant à
$\symup{\underline{\operatorname{Pic}}_{X/k}}$.

\result{\uline{Remarque}}{6.7} Soit plus généralement $\symup{f:X'\to X}$ un
morphisme surjectif de préschémas propres sur $\symup{k}$. Alors des considérations de
descente non plate conduisent à conjecturer que
$\symup{\underline{\operatorname{Pic}}_{X/k}\to\underline{\operatorname{Pic}}_{X'/k}}$ est un morphisme \uline{affine}, ce qui impliquerait en particulier (en divisant par les
composantes connexes des éléments neutres) que l'homomorphisme correspondant
sur les groupes de Néron-Severi est injectif modulo torsion. C'est ce qu'on
peut vérifier par la théorie des intersections lorsque $\symup{X}$ et $\symup{X'}$ sont non
singuliers. La réponse ne semble pas connue dans tout autre cas.

\result{\uline{Remarque}}{6.8} Contrairement à ce qu'on pourrait penser, la
considération des schémas de Picard de schémas algébriques à éléments
nilpotents est utile, et même indispensable, dans diverses questions. Ainsi,
lorsque $\symup{X}$ est un schéma projectif, disons régulier, et $\symup{Y}$ une
section hyperplane, il y a lieu de considérer les ``voisinages
infinitésimaux'' de tous ordres $\symup{X_n}$ de $\symup{Y}$, et les schémas de Picard
$\symup{\underline{\operatorname{Pic}}_{X_n/k}}$ ; lorsque $\symup{X}$ est irréductible de dimension
$\symup{\geqslant4}$ (resp. $\symup{\geqslant3}$), le morphisme canonique
\[
\symup{\underline{\operatorname{Pic}}_{X/k}\to\underline{\operatorname{Pic}}_{X_n/k}
\qquad(n\text{ grand})}
\]
est un isomorphisme (resp. induit un isomorphisme pour les images inverses des
sous-groupes de torsion de Néron-Severi), résultat qui nous sera utile dans
l'étude qualitative des schémas de Picard dans l'exposé suivant. De même, la
considération des schémas de Picard de certaines courbes à éléments nilpotents
et les théorèmes fondamentaux de géométrie formelle [3], permettent de
vérifier, dans le cas d'égales caractéristiques, une conjecture de MUMFORD,
savoir que pour tout anneau local noethérien normal complet $\symup{A}$ de dimension
2, le groupe des classes de diviseurs de $\symup{A}$ peut être considéré comme
l'ensemble des points rationnels sur $\symup{k}$ d'un

\src{161}
groupe algébrique $\symup{G}$ sur le corps résiduel $\symup{k}$ ($\symup{G}$ étant d'ailleurs
déterminé canoniquement une fois qu'on se donne un corps de représentants dans
$\symup{A}$). Dans le cas où $\symup{A}$ est de dimension quelconque, il est plausible qu'il
existe un pro-groupe algébrique sur $\symup{k}$ jouant le même rôle que $\symup{G}$ ci-dessus,
qui se construit, dans le cas où on peut ``désingulariser''
$\symup{\operatorname{Spec}(A)}$, comme une limite projective de schémas de Picard de
schémas projectifs (à éléments nilpotents) convenables sur $\symup{k}$.

\begin{center}BIBLIOGRAPHIE\end{center}
\addcontentsline{toc}{section}{Bibliographie}

\begin{description}
\item[{\normalfont[1]}] CHEVALLEY (Claude). - Sur la théorie de Picard,
Amer. J. of Math., t. 82, 1960, p. 435-490.

\item[{\normalfont[2]}] GROTHENDIECK (A.) et DIEUDONNÉ (J.). - Éléments de géométrie
algébrique, Chapitre I et suivants. - Paris, Presses universitaires de France,
1960-1961 (Institut des Hautes Etudes Scientifiques, Publications
mathématiques 4, 8, \ldots).

\item[{\normalfont[3]}] GROTHENDIECK (Alexander). - Géométrie formelle et géométrie
algébrique, Séminaire Bourbaki, t. 11, 1958/59,
n\textsuperscript{o} 182, 28 p.

\item[{\normalfont[4]}] GROTHENDIECK (Alexander). - Techniques de descente et théorèmes
d'existence en géométrie algébrique, I-IV, Séminaire Bourbaki, t. 12,
1959/60, n\textsuperscript{o} 190, 29 p. et n\textsuperscript{o} 195, 22 p. ;
t. 13, 1960/61, n\textsuperscript{o} 212, 20 p. et
n\textsuperscript{o} 221, 28 p.

\item[{\normalfont[5]}] GROTHENDIECK (Alexander). - Techniques de construction en
géométrie analytique, I-X, Séminaire Cartan, t. 13, 1960/61 :
Déformation des structures analytiques complexes, n\textsuperscript{o} 7-17.

\item[{\normalfont[6]}] GROTHENDIECK (Alexander). - Séminaire de géométrie algébrique,
rédigé par Lichtenbaum, Harvard University, 1961 (à paraître).

\item[{\normalfont[7]}] LANG (Serge). - Abelian varieties. - New York,
Interscience Publishers, 1959 (Interscience Tracts in pure and applied
Mathematics, 7).

\item[{\normalfont[8]}] MUMFORD (D.). - An elementary theorem in geometric invariant
theory, Bull. Amer. math. Soc., t. 67, 1961, p. 483-486.

\item[{\normalfont[9]}] MUMFORD (D.). - The topology of normal singularities of an
algebraic surface and a criterion for simplicity. - Paris, Presses
universitaires de France, 1961 (Institut des Hautes Etudes Scientifiques.
Publications mathématiques, 9).

\item[{\normalfont[10]}] MUMFORD (D.) and TATE (J.). - Séminaire de géométrie
algébrique, Harvard University, Spring Term 1962 (à paraître).

\item[{\normalfont[11]}] OORT (F.). - Sur le schéma de Picard,
Bull. Soc. math. France, t. 90, 1962 (à paraître).

\item[{\normalfont[12]}] SESHADRI (C. S.). - Thèse à paraître dans
Annali di Matematica.
\end{description}

\begin{center}
  \rule{.20\linewidth}{.3pt}
\end{center}
